\hypertarget{huawei-ai-certification-ai-final-exam-lab}{%
\section{Huawei AI Certification: AI Final Exam
Lab}\label{huawei-ai-certification-ai-final-exam-lab}}

\hypertarget{handwritten-digit-recognition-mnist}{%
\subsection{Handwritten Digit Recognition
(MNIST)}\label{handwritten-digit-recognition-mnist}}

\textbf{Objective:} Master the basic process of TensorFlow computing and
build neural networks to classify MNIST hand-written digits. We will
evaluate a Deep Neural Network (DNN) against a Convolutional Neural
Network (CNN) architecture.

\begin{lstlisting}[language=Python]
import os
import numpy as np
import matplotlib.pyplot as plt

import tensorflow as tf
from tensorflow import keras
from tensorflow.keras import datasets, layers, optimizers
from tensorflow.keras.layers import Conv2D, MaxPool2D, Dropout, Flatten, Dense
from tensorflow.keras.models import Sequential, load_model

# Ensure inline plotting
%matplotlib inline
\end{lstlisting}

\begin{lstlisting}
WARNING: All log messages before absl::InitializeLog() is called are written to STDERR
I0000 00:00:1777685359.528055 2702417 cudart_stub.cc:31] Could not find cuda drivers on your machine, GPU will not be used.
I0000 00:00:1777685359.587772 2702417 cpu_feature_guard.cc:227] This TensorFlow binary is optimized to use available CPU instructions in performance-critical operations.
To enable the following instructions: AVX2 FMA, in other operations, rebuild TensorFlow with the appropriate compiler flags.
WARNING: All log messages before absl::InitializeLog() is called are written to STDERR
I0000 00:00:1777685361.442989 2702417 cudart_stub.cc:31] Could not find cuda drivers on your machine, GPU will not be used.
\end{lstlisting}

\hypertarget{dataset-acquisition}{%
\subsubsection{1. Dataset Acquisition}\label{dataset-acquisition}}

Load the MNIST dataset natively into training and test variables, and
construct one-hot encoded tensors for classification.

\begin{lstlisting}[language=Python]
(x_train, y_train), (x_test, y_test) = datasets.mnist.load_data()
print("First 5 labels:", y_train[:5])
print("x_train shape:", x_train.shape, "y_train shape:", y_train.shape)
print("x_test shape :", x_test.shape, "y_test shape :", y_test.shape)

n_classes = 10
y_train_oh = keras.utils.to_categorical(y_train, n_classes)
y_test_oh = keras.utils.to_categorical(y_test, n_classes)
print("One-hot sample:\n", y_train_oh[:5])
\end{lstlisting}

\begin{lstlisting}
First 5 labels: [5 0 4 1 9]
x_train shape: (60000, 28, 28) y_train shape: (60000,)
x_test shape : (10000, 28, 28) y_test shape : (10000,)
One-hot sample:
 [[0. 0. 0. 0. 0. 1. 0. 0. 0. 0.]
 [1. 0. 0. 0. 0. 0. 0. 0. 0. 0.]
 [0. 0. 0. 0. 1. 0. 0. 0. 0. 0.]
 [0. 1. 0. 0. 0. 0. 0. 0. 0. 0.]
 [0. 0. 0. 0. 0. 0. 0. 0. 0. 1.]]
\end{lstlisting}

\hypertarget{data-visualization}{%
\subsubsection{2. Data Visualization}\label{data-visualization}}

Extract and plot exactly how the computer natively parses the numeric
pixel grids.

\begin{lstlisting}[language=Python]
plt.figure(figsize=(6, 6))
for i in range(25):
    plt.subplot(5, 5, i + 1)
    plt.imshow(x_train[i], cmap="gray")
    plt.axis("off")
plt.suptitle("First 25 MNIST training images")
plt.tight_layout()
plt.show()
\end{lstlisting}

\begin{figure}[H]
\centering
\includegraphics[width=0.85\linewidth]{"AI Final Exam Lab/AI Final Exam Lab_files/AI Final Exam Lab_5_0.png"}
\end{figure}

\hypertarget{deep-neural-network-dnn-construction}{%
\subsubsection{3. Deep Neural Network (DNN)
Construction}\label{deep-neural-network-dnn-construction}}

We flat-pack the 28x28 matrices into 1D (784) tensors, and map them
linearly across multiple dense feed-forward node layers.

\begin{lstlisting}[language=Python]
# Preprocess for DNN (flatten + normalize)
X_train_dnn = x_train.reshape(60000, 784).astype("float32") / 255.0
X_test_dnn = x_test.reshape(10000, 784).astype("float32") / 255.0

dnn = Sequential([
    Dense(512, activation="relu", input_dim=784),
    Dense(256, activation="relu"),
    Dense(128, activation="relu"),
    Dense(n_classes, activation="softmax"),
])
dnn.summary()

dnn.compile(
    loss=keras.losses.categorical_crossentropy,
    optimizer=optimizers.Adam(0.001),
    metrics=["accuracy"],
)

dnn_history = dnn.fit(
    X_train_dnn, y_train_oh,
    batch_size=128, epochs=10, verbose=2,
    validation_data=(X_test_dnn, y_test_oh),
)

dnn_score = dnn.evaluate(X_test_dnn, y_test_oh, verbose=0)
print("DNN Test loss:", dnn_score[0])
print("DNN Test accuracy:", dnn_score[1])
\end{lstlisting}

\begin{lstlisting}
/home/ansinitro/AITU/ai-nn/midterm/.venv/lib/python3.12/site-packages/keras/src/layers/core/dense.py:107: UserWarning: Do not pass an `input_shape`/`input_dim` argument to a layer. When using Sequential models, prefer using an `Input(shape)` object as the first layer in the model instead.
  super().__init__(activity_regularizer=activity_regularizer, **kwargs)
E0000 00:00:1777685363.240408 2702417 cuda_executor.cc:1737] INTERNAL: CUDA Runtime error: Failed call to cudaGetRuntimeVersion: Error loading CUDA libraries. GPU will not be used.: Error loading CUDA libraries. GPU will not be used.
W0000 00:00:1777685363.240779 2702502 cuda_executor.cc:1755] Failed to determine cuDNN version (Note that this is expected if the application doesn't link the cuDNN plugin): INTERNAL: cuDNN error: CUDNN_STATUS_INTERNAL_ERROR
W0000 00:00:1777685363.255852 2702417 gpu_device.cc:2365] Cannot dlopen some GPU libraries. Please make sure the missing libraries mentioned above are installed properly if you would like to use GPU. Follow the guide at https://www.tensorflow.org/install/gpu for how to download and setup the required libraries for your platform.
Skipping registering GPU devices...
\end{lstlisting}

\begin{lstlisting}
Epoch 1/10


W0000 00:00:1777685363.614474 2702417 cpu_allocator_impl.cc:82] Allocation of 188160000 exceeds 10% of free system memory.


469/469 - 4s - 10ms/step - accuracy: 0.9319 - loss: 0.2288 - val_accuracy: 0.9677 - val_loss: 0.1055
Epoch 2/10
469/469 - 3s - 6ms/step - accuracy: 0.9737 - loss: 0.0863 - val_accuracy: 0.9723 - val_loss: 0.0853
Epoch 3/10
469/469 - 3s - 6ms/step - accuracy: 0.9827 - loss: 0.0555 - val_accuracy: 0.9770 - val_loss: 0.0739
Epoch 4/10
469/469 - 3s - 6ms/step - accuracy: 0.9876 - loss: 0.0391 - val_accuracy: 0.9724 - val_loss: 0.0953
Epoch 5/10
469/469 - 3s - 6ms/step - accuracy: 0.9893 - loss: 0.0317 - val_accuracy: 0.9754 - val_loss: 0.0865
Epoch 6/10
469/469 - 3s - 7ms/step - accuracy: 0.9912 - loss: 0.0272 - val_accuracy: 0.9789 - val_loss: 0.0750
Epoch 7/10
469/469 - 3s - 7ms/step - accuracy: 0.9934 - loss: 0.0205 - val_accuracy: 0.9789 - val_loss: 0.0769
Epoch 8/10
469/469 - 4s - 8ms/step - accuracy: 0.9928 - loss: 0.0217 - val_accuracy: 0.9798 - val_loss: 0.0792
Epoch 9/10
469/469 - 3s - 7ms/step - accuracy: 0.9940 - loss: 0.0179 - val_accuracy: 0.9780 - val_loss: 0.0850
Epoch 10/10
469/469 - 3s - 7ms/step - accuracy: 0.9953 - loss: 0.0142 - val_accuracy: 0.9799 - val_loss: 0.0817
DNN Test loss: 0.08167005330324173
DNN Test accuracy: 0.9799000024795532
\end{lstlisting}

\hypertarget{convolutional-neural-network-cnn-construction}{%
\subsubsection{4. Convolutional Neural Network (CNN)
Construction}\label{convolutional-neural-network-cnn-construction}}

We preserve the spatial 2D matrix shape, allowing the CNN kernels to
systematically extract positional hierarchies and edge gradients
natively, maximizing feature recognition accuracy!

\begin{lstlisting}[language=Python]
X_train_cnn = x_train.reshape(-1, 28, 28, 1).astype("float32") / 255.0
X_test_cnn = x_test.reshape(-1, 28, 28, 1).astype("float32") / 255.0

cnn = Sequential()
cnn.add(Conv2D(32, kernel_size=(3, 3), padding="same", activation="relu", input_shape=(28, 28, 1)))
cnn.add(MaxPool2D(pool_size=(2, 2)))
cnn.add(Conv2D(64, kernel_size=(3, 3), padding="valid", activation="relu"))
cnn.add(MaxPool2D(pool_size=(2, 2)))
cnn.add(Dropout(0.25))
cnn.add(Flatten())
cnn.add(Dense(128, activation="relu"))
cnn.add(Dropout(0.5))
cnn.add(Dense(n_classes, activation="softmax"))
cnn.summary()

cnn.compile(
    loss="categorical_crossentropy",
    optimizer="adam",
    metrics=["accuracy"],
)

cnn_history = cnn.fit(
    X_train_cnn, y_train_oh,
    epochs=5, batch_size=128, verbose=2,
    validation_data=(X_test_cnn, y_test_oh),
)

cnn_score = cnn.evaluate(X_test_cnn, y_test_oh, verbose=0)
print("CNN Test Loss:", cnn_score[0])
print("CNN Test Accuracy:", cnn_score[1])
\end{lstlisting}

\begin{lstlisting}
/home/ansinitro/AITU/ai-nn/midterm/.venv/lib/python3.12/site-packages/keras/src/layers/convolutional/base_conv.py:113: UserWarning: Do not pass an `input_shape`/`input_dim` argument to a layer. When using Sequential models, prefer using an `Input(shape)` object as the first layer in the model instead.
  super().__init__(activity_regularizer=activity_regularizer, **kwargs)
\end{lstlisting}

\begin{lstlisting}
Epoch 1/5


W0000 00:00:1777685397.829101 2702417 cpu_allocator_impl.cc:82] Allocation of 188160000 exceeds 10% of free system memory.


469/469 - 21s - 45ms/step - accuracy: 0.9106 - loss: 0.2942 - val_accuracy: 0.9825 - val_loss: 0.0550
Epoch 2/5
469/469 - 17s - 37ms/step - accuracy: 0.9717 - loss: 0.0955 - val_accuracy: 0.9878 - val_loss: 0.0369
Epoch 3/5
469/469 - 18s - 38ms/step - accuracy: 0.9784 - loss: 0.0724 - val_accuracy: 0.9897 - val_loss: 0.0321
Epoch 4/5
469/469 - 16s - 35ms/step - accuracy: 0.9821 - loss: 0.0607 - val_accuracy: 0.9892 - val_loss: 0.0312
Epoch 5/5
469/469 - 16s - 34ms/step - accuracy: 0.9844 - loss: 0.0509 - val_accuracy: 0.9915 - val_loss: 0.0260
CNN Test Loss: 0.026009762659668922
CNN Test Accuracy: 0.9915000200271606
\end{lstlisting}

\hypertarget{architectural-evaluation-and-training-validation-curves}{%
\subsubsection{5. Architectural Evaluation and Training Validation
Curves}\label{architectural-evaluation-and-training-validation-curves}}

\begin{lstlisting}[language=Python]
def plot_history(history, title):
    fig, axes = plt.subplots(1, 2, figsize=(11, 4))
    axes[0].plot(history.history["loss"], "o-", label="train")
    axes[0].plot(history.history["val_loss"], "s-", label="val")
    axes[0].set_title(f"{title} — Loss")
    axes[0].set_xlabel("epoch"); axes[0].set_ylabel("loss")
    axes[0].legend(); axes[0].grid(alpha=0.3)

    axes[1].plot(history.history["accuracy"], "o-", label="train")
    axes[1].plot(history.history["val_accuracy"], "s-", label="val")
    axes[1].set_title(f"{title} — Accuracy")
    axes[1].set_xlabel("epoch"); axes[1].set_ylabel("accuracy")
    axes[1].legend(); axes[1].grid(alpha=0.3)

    plt.tight_layout()
    plt.show()

plot_history(dnn_history, "DNN")
plot_history(cnn_history, "CNN")

# ---------- Model comparison bar chart ----------
fig, ax = plt.subplots(figsize=(6, 4))
names = ["DNN", "CNN"]
losses = [dnn_score[0], cnn_score[0]]
accs = [dnn_score[1], cnn_score[1]]
xs = np.arange(len(names))
w = 0.35
b1 = ax.bar(xs - w / 2, losses, w, label="test loss", color="#e76f51")
b2 = ax.bar(xs + w / 2, accs, w, label="test accuracy", color="#2a9d8f")
ax.set_xticks(xs); ax.set_xticklabels(names)
ax.set_title("DNN vs CNN — test metrics")
ax.legend(); ax.grid(axis="y", alpha=0.3)
for b in (*b1, *b2):
    ax.text(b.get_x() + b.get_width() / 2, b.get_height(), f"{b.get_height():.4f}", ha="center", va="bottom", fontsize=9)
plt.tight_layout()
plt.show()
\end{lstlisting}

\begin{figure}[H]
\centering
\includegraphics[width=0.85\linewidth]{"AI Final Exam Lab/AI Final Exam Lab_files/AI Final Exam Lab_11_0.png"}
\end{figure}

\begin{figure}[H]
\centering
\includegraphics[width=0.85\linewidth]{"AI Final Exam Lab/AI Final Exam Lab_files/AI Final Exam Lab_11_1.png"}
\end{figure}

\begin{figure}[H]
\centering
\includegraphics[width=0.85\linewidth]{"AI Final Exam Lab/AI Final Exam Lab_files/AI Final Exam Lab_11_2.png"}
\end{figure}

\hypertarget{confusion-matrices-and-predictions}{%
\subsubsection{6. Confusion Matrices and
Predictions}\label{confusion-matrices-and-predictions}}

Finally, evaluating the exact false-positive failure modes between
architectures natively.

\begin{lstlisting}[language=Python]
def plot_confusion(model, X, y_true, title):
    probs = model.predict(X, verbose=0)
    y_pred = np.argmax(probs, axis=1)
    cm = np.zeros((n_classes, n_classes), dtype=int)
    for t, p in zip(y_true, y_pred):
        cm[t, p] += 1

    fig, ax = plt.subplots(figsize=(6, 5))
    im = ax.imshow(cm, cmap="Blues")
    ax.set_title(f"Confusion matrix — {title}")
    ax.set_xlabel("predicted"); ax.set_ylabel("true")
    ax.set_xticks(range(n_classes)); ax.set_yticks(range(n_classes))
    fig.colorbar(im, ax=ax)
    thresh = cm.max() / 2.0
    for i in range(n_classes):
        for j in range(n_classes):
            ax.text(j, i, cm[i, j], ha="center", va="center", fontsize=8, color="white" if cm[i, j] > thresh else "black")
    plt.tight_layout()
    plt.show()
    return y_pred

dnn_pred = plot_confusion(dnn, X_test_dnn, y_test, "DNN")
cnn_pred = plot_confusion(cnn, X_test_cnn, y_test, "CNN")
\end{lstlisting}

\begin{figure}[H]
\centering
\includegraphics[width=0.85\linewidth]{"AI Final Exam Lab/AI Final Exam Lab_files/AI Final Exam Lab_13_0.png"}
\end{figure}

\begin{figure}[H]
\centering
\includegraphics[width=0.85\linewidth]{"AI Final Exam Lab/AI Final Exam Lab_files/AI Final Exam Lab_13_1.png"}
\end{figure}

\hypertarget{model-deployment-simulation}{%
\subsubsection{Model Deployment
Simulation}\label{model-deployment-simulation}}

Running a mock inference pipeline on 20 raw tensor inputs to determine
deployment readiness.

\begin{lstlisting}[language=Python]
def res_visual(n):
    probs = cnn.predict(X_test_cnn[0:n], verbose=0)
    final_opt_a = np.argmax(probs, axis=1)

    fig, ax = plt.subplots(nrows=int(n / 5), ncols=5, figsize=(8, 1.8 * (n // 5)))
    ax = ax.flatten()
    print(f"Prediction results of the first {n} images:")
    for i in range(n):
        img = X_test_cnn[i].reshape((28, 28))
        ax[i].imshow(img, cmap="Greys")
        ok = final_opt_a[i] == y_test[i]
        ax[i].set_title(f"pred {final_opt_a[i]} / true {y_test[i]}", fontsize=8, color="green" if ok else "red")
        ax[i].axis("off")
    plt.tight_layout()
    plt.show()

res_visual(20)
\end{lstlisting}

\begin{lstlisting}
Prediction results of the first 20 images:
\end{lstlisting}

\begin{figure}[H]
\centering
\includegraphics[width=0.85\linewidth]{"AI Final Exam Lab/AI Final Exam Lab_files/AI Final Exam Lab_15_1.png"}
\end{figure}
